\documentclass[11pt,a4paper]{article}
\usepackage[utf8]{inputenc}
\usepackage[T1]{fontenc}
\usepackage[french]{babel}
\usepackage{graphicx}
\usepackage{amsmath}
\usepackage{booktabs}
\usepackage{hyperref}
\usepackage[margin=2.5cm]{geometry}

\graphicspath{{figures/}}

\title{Micro-réseau intelligent : prévision de la consommation et de la production solaire, tarification dynamique}
\author{Dev A \and Dev B \and Dev C}
\date{\today}

\begin{document}
\maketitle

\begin{abstract}
% Résumé du projet en quelques lignes.
\end{abstract}

\tableofcontents

\section{Introduction}
% Contexte des smart grids, objectifs, contributions.

\section{Données}
\subsection{London Smart Meters}
% Description, période, nombre de foyers, nettoyage.
\subsection{Irradiance solaire}
% Source, variables, conversion en production PV.

\section{Méthodologie}
\subsection{Prévision de consommation : régression linéaire (Dev A)}
% Variables calendaires, retards, régularisation.
\subsection{Prévision solaire : LSTM (Dev B)}
% Architecture, fenêtre temporelle, entraînement.
\subsection{Agent de tarification : réseau dense (Dev C)}
% État, sortie, fonction cible, entraînement.
\subsection{Environnement de simulation}
% Bilan énergétique, batterie, élasticité prix.

\section{Résultats}
\begin{table}[h]
\centering
\begin{tabular}{lcccc}
\toprule
Modèle & MAE & RMSE & MAPE (\%) & $R^2$ \\
\midrule
Régression linéaire & & & & \\
LSTM solaire & & & & \\
\bottomrule
\end{tabular}
\caption{Performances des modèles de prévision sur le jeu de test.}
\end{table}

% \begin{figure}[h]
%   \centering
%   \includegraphics[width=\textwidth]{simulation.png}
%   \caption{Simulation du micro-réseau.}
% \end{figure}

\section{Discussion}
% Limites, pistes d'amélioration.

\section{Conclusion}

\end{document}
